%------------------------------------------------------------------------------%

\usepackage{integracao_e_series}

\title{Teorema do Valor Médio para Integrais}

%------------------------------------------------------------------------------%
\begin{document}

\addtitle

%------------------------------------------------------------------------------%
\section{Recapitulando}

%------------------------------------------------------------------------------%
\begin{frame}{Derivadas e Integrais}
\emph{Derivada}
\tabto{37mm}Taxa de variação de $f$
\[
  \hspace{45mm}
  f'(x) = \frac{df}{dx}
\]

\pause
\emph{Integral indefinida}
\tabto{37mm}Antiderivada de $f$
\[
  \hspace{45mm}
  F(x) = \int f(x)\,dx \qquad\Rightarrow\qquad \frac{dF}{dx} = f(x)
\]

\pause
\emph{Integral definida}
\tabto{37mm}Diferença de áreas sob o gráfico de $f$
\[
  \hspace{45mm}
  \text{``Área''} = \int_a^b f(x)\,dx
\]
\end{frame}

%------------------------------------------------------------------------------%
\section{Teorema do Valor Médio para Integrais}

%------------------------------------------------------------------------------%
\begin{frame}{Teorema do Valor Médio para Integrais}
  Seja $f$ uma função {\color{blue} contínua} em $[a, b]$, \pause
  então existe ${\color{blue} c} \in (a, b)$ \pause
  tal que
  \[
    f({\color{blue}c})(b-a)
    \;=\; \int_a^b f(x)\,dx
  \]
  \pause
  ou
  \[
    f({\color{blue}c})
    \;=\; \frac{1}{b-a}\int_a^b f(x)\,dx
  \]
\end{frame}

%------------------------------------------------------------------------------%
\begin{frame}{Valor Médio}
  ${\color{blue}f(c)}$ é o \emph{Valor Médio} de $f$ em $[a, b]$
  que denotamos por ${\color{blue}f_{\text{médio}}}$.
\end{frame}

%------------------------------------------------------------------------------%
\begin{frame}{Valor Médio}
  \centering
  \vspace{\baselineskip}
  \input{fig_valor_medio_media}
\end{frame}

%------------------------------------------------------------------------------%
\section{Demonstração}

%------------------------------------------------------------------------------%
\begin{frame}{Demonstração}
  Queremos mostrar que existe $c \in (a, b)$ tal que
  \[
    f(c) = \frac{1}{b-a}\int_a^b f(x)\,dx
  \]
\end{frame}

%------------------------------------------------------------------------------%
\begin{frame}{Demonstração}
  \centering
  \vspace{\baselineskip}
  \input{fig_valor_medio_media}
\end{frame}

%------------------------------------------------------------------------------%
\begin{frame}{Demonstração}
\onslide<+->{Teorema de Weierstrass \\
\qquad $f$ é \emph{contínua} em $[a, b]$
então existe um \emph{mínimo} $f(u)$ e um \emph{máximo} $f(v)$}

\vspace{0.5\baselineskip}
\onslide<+->{Como a integral \emph{preserva a desigualdade}}
\begin{align*}
  \onslide<+->{\int_a^b f(u)\,dx \;\leq\; & \int_a^b f(x)\,dx \;\leq\; \int_a^b f(v) \,dx \\[3mm]}
  \onslide<+->{f(u)(b-a) \;\leq\;         & \int_a^b f(x)\, dx \;\leq \;f(v)(b-a)         \\[3mm]}
  \onslide<+->{f(u) \;\leq\;               \frac{1}{b-a} & \int_a^b f(x)\,dx \;\leq\; f(v)}
\end{align*}
\end{frame}

%------------------------------------------------------------------------------%
\begin{frame}{Demonstração}
Valor médio
\[
  \fmedio = \frac{1}{b-a} \int_a^b f(x)\,dx
\]
\pause
\[
  f(u) \;\leq\; f_{\text{médio}} \;\leq\; f(v)
\]

\pause
\vspace{\baselineskip}
Teorema do Valor Intermediário \\
\qquad $f$ precisa assumir o valor \(\fmedio\) em algum $c\in[a, b]$
\end{frame}

%------------------------------------------------------------------------------%
\addtocounter{exemplo}{1}
\begin{frame}{Exemplo \theexemplo}
  Cálcule \(\)
\end{frame}

\begin{frame}{Exemplo \theexemplo}
  \begin{align*}
  \end{align*}
\end{frame}

%------------------------------------------------------------------------------%
\section{Lista Mínima}

%------------------------------------------------------------------------------%
\begin{frame}{Lista Mínima}
  Estudar a Seção 2.? da Apostila

  Exercícios:

  \vfill
  Atenção: A prova é baseada no livro, não nas apresentações
\end{frame}

%------------------------------------------------------------------------------%
\end{document}
%------------------------------------------------------------------------------%
